\documentclass[11pt]{article}

\usepackage[a4paper,margin=1in]{geometry}
\usepackage{amsmath,amssymb}
\usepackage{booktabs}

\title{Reaction--Diffusion Model with a Periodic Gaussian Source}
\author{}
\date{}

\newcommand{\dd}{\,\mathrm{d}}
\newcommand{\R}{\mathbb{R}}

\begin{document}

\maketitle

\section{Spatial domain}

The computational domain is a unit disk with a circular hole,
\begin{equation}
  \Omega
  = \left\{\boldsymbol{x}=(x,y)\in\R^2:
      \lVert\boldsymbol{x}\rVert < 1\right\}
    \setminus
    \left\{\boldsymbol{x}\in\R^2:
      \lVert\boldsymbol{x}-\boldsymbol{x}_{\mathrm{hole}}\rVert
      \leq r_{\mathrm{hole}}\right\},
\end{equation}
where
\(
  \boldsymbol{x}_{\mathrm{hole}}=(-0.25,0)
\)
and
\(
  r_{\mathrm{hole}}=0.3
\).

\section{PDE model}

For final time \(T=5\), find the scalar field
\(u:\Omega\times[0,T]\to\R\) satisfying
\begin{align}
  \frac{\partial u}{\partial t}
  - \nabla\!\cdot\!\left(D\nabla u\right)
  + r u
  &= F(\boldsymbol{x},t)
  && \text{in }\Omega\times(0,T],
  \label{eq:pde}\\
  \boldsymbol{n}\cdot D\nabla u
  &=0
  && \text{on }\partial\Omega\times(0,T],
  \label{eq:neumann}\\
  u(\boldsymbol{x},0)
  &=u_0(\boldsymbol{x})
  && \text{in }\Omega.
  \label{eq:initial}
\end{align}
Here \(\boldsymbol{n}\) denotes the outward unit normal on both the outer
boundary and the boundary of the hole.  The diffusion tensor and reaction
coefficient are
\begin{equation}
  D=\operatorname{diag}(a_1,a_2),
  \qquad a_1=0.8,
  \qquad a_2=0.6,
  \qquad r=1.
\end{equation}

The initial depolarization is the Gaussian
\begin{equation}
  u_0(\boldsymbol{x})
  = A_0
    \exp\!\left(
      -\frac{\lVert\boldsymbol{x}-\boldsymbol{x}_0\rVert^2}
             {2\sigma_0^2}
    \right),
  \qquad
  A_0=1,
  \quad
  \boldsymbol{x}_0=(0.4,0),
  \quad
  \sigma_0=0.12.
\end{equation}

\subsection{Periodic external source}

The source is separable,
\begin{equation}
  F(\boldsymbol{x},t)=f(\boldsymbol{x})g(t),
\end{equation}
with spatial profile
\begin{equation}
  f(\boldsymbol{x})
  = A_F
    \exp\!\left(
      -\frac{\lVert\boldsymbol{x}-\boldsymbol{x}_F\rVert^2}
             {2\sigma_x^2}
    \right)
\end{equation}
and smooth periodic heartbeat
\begin{equation}
  g(t)
  = \exp\!\left[
      -\frac{
        \sin^2\!\left(\pi(t-\phi)/P\right)
      }{
        2\left(\pi\sigma_t/P\right)^2
      }
    \right].
  \label{eq:heartbeat}
\end{equation}
The function \(g\) is exactly \(P\)-periodic and has unit-height peaks at
\(t=\phi+kP\), \(k\in\mathbb{Z}\).  With \(P=1\), one heartbeat occurs every
second.

\section{Model and discretization parameters}

\begin{center}
\begin{tabular}{@{}lll@{}}
\toprule
Quantity & Symbol or code name & Value \\
\midrule
Final time & \(T\) & \(5\) \\
Number of time steps & \texttt{num\_steps} & \(500\) \\
Time-step size & \(\Delta t=T/\texttt{num\_steps}\) & \(0.01\) \\
Theta-scheme parameter & \(\theta\) & \(1\) (backward Euler) \\
Target mesh size & \(h\) & \(0.025\) \\
Finite-element degree & \(p\) & \(2\) \\
Diffusion coefficients & \(a_1,a_2\) & \(0.8,0.6\) \\
Reaction coefficient & \(r\) & \(1\) \\
Initial amplitude & \(A_0\) & \(1\) \\
Initial center & \(\boldsymbol{x}_0\) & \((0.4,0)\) \\
Initial spatial width & \(\sigma_0\) & \(0.12\) \\
Source amplitude & \(A_F\) & \(1\) \\
Source center & \(\boldsymbol{x}_F\) & \((0.4,0)\) \\
Source spatial width & \(\sigma_x\) & \(0.12\) \\
Heartbeat period & \(P\) & \(1\) \\
Heartbeat time width & \(\sigma_t\) & \(0.05\) \\
Heartbeat phase & \(\phi\) & \(0\) \\
Hole center & \(\boldsymbol{x}_{\mathrm{hole}}\) & \((-0.25,0)\) \\
Hole radius & \(r_{\mathrm{hole}}\) & \(0.3\) \\
\bottomrule
\end{tabular}
\end{center}

\section{Finite-element and time discretization}

Let \(\mathcal{T}_h\) be a triangular mesh of \(\Omega\), and define the
continuous quadratic finite-element space
\begin{equation}
  V_h
  = \left\{v_h\in C^0(\overline{\Omega}):
      v_h|_K\in\mathbb{P}_2(K)
      \text{ for every }K\in\mathcal{T}_h\right\}.
\end{equation}
No essential boundary conditions are imposed because the zero-flux condition
in~\eqref{eq:neumann} is natural.  Introduce the bilinear form
\begin{equation}
  B(w,v)
  = \int_\Omega (D\nabla w)\cdot\nabla v\dd\boldsymbol{x}
    + \int_\Omega r wv\dd\boldsymbol{x}.
\end{equation}

Set \(t_n=n\Delta t\) and let \(u_h^0\in V_h\) be the finite-element
interpolant of \(u_0\).  Given \(u_h^n\), the theta scheme computes
\(u_h^{n+1}\in V_h\) such that, for every \(v_h\in V_h\),
\begin{align}
  (u_h^{n+1},v_h)
  + \Delta t\,\theta B(u_h^{n+1},v_h)
  ={}&(u_h^n,v_h)
  - \Delta t\,(1-\theta)B(u_h^n,v_h)
  \nonumber\\
  &+\Delta t\left(
      \theta F(\boldsymbol{x},t_{n+1})
      +(1-\theta)F(\boldsymbol{x},t_n),v_h
    \right),
  \label{eq:theta-scheme}
\end{align}
where \((w,v)=\int_\Omega wv\dd\boldsymbol{x}\).  Thus, the source is
evaluated with the same theta weighting as the spatial operator.  The current
choice \(\theta=1\) reduces~\eqref{eq:theta-scheme} to backward Euler; choosing
\(\theta=1/2\) gives Crank--Nicolson.

For a basis \(\{\varphi_i\}\) of \(V_h\), define
\begin{align}
  M_{ij} &= (\varphi_j,\varphi_i),\\
  K_{ij} &= B(\varphi_j,\varphi_i),\\
  b_{F,i} &= (f,\varphi_i).
\end{align}
The linear system at time step \(n+1\) is
\begin{equation}
  \left[M+\Delta t\,\theta K\right]U^{n+1}
  = \left[M-\Delta t(1-\theta)K\right]U^n
    +\Delta t\left[
      \theta g(t_{n+1})+(1-\theta)g(t_n)
    \right]b_F.
\end{equation}
Because \(D\), \(r\), \(\Delta t\), and \(\theta\) are constant, the
left-hand matrix is assembled and factorized once, then reused for all time
steps.  The implementation uses a direct LU solve through PETSc.

\end{document}
